\section*{4 Factors of Arithmetic Branching Programs}
In this section, we prove that ABPs are closed under factoring.
\setcounter{nttheorem}{10}
We prove Theorem 10 in the rest of this section. First observe that it suffices to prove the $\operatorname{poly}(s)$ size upper bound for the irreducible factors of $p$. This yields the same bound for all the factors. The case when $p=q^e$ is proved in Section 3.2. So it remains to consider the general case when $p=p_1^{e_1}\cdots p_m^{e_m}$, for $m\geq2$, where $p_1,\ldots,p_m$ are the different irreducible factors of $p$. We want to prove an ABP size upper bound for an irreducible factor, say $p_1$.
We start by several transformations on the input polynomial $p(\mathbf z)$, where $\mathbf z=(z_1,\ldots,z_n)$.
\begin{enumerate}
\item As described in Section 3.3, taking $k=e_1-1$ times the derivative w.r.t. some variable, say $z_1$, we get the polynomial $p'(\mathbf z)=\frac{\partial^kp(\mathbf z)}{\partial z_1^k}$, where the factor $p_1$ has multiplicity1.
\item Next, by Lemma 7, we transform $p'(\mathbf z)$ to a polynomial $p''(x,\mathbf z)$ that is monic in $x$, for a new variable $x$. Thereby also the factors of $p'(\mathbf z)$ are transformed, maintaining their irreducibilty and multiplicity. The degree of $p''$ is twice the degree of $p'$.
\item At this point, we may have to shift the variables $\mathbf z$ as described in Section 3.5 to ensure the properties needed for starting the Hensel lifting. This shift preserves the monicness and the irreducibility of the factors.
\item Finally, the transformation to a bivariate polynomial is explained in Section 3.6. This yields polynomial $p'''(x,y,\mathbf z)$, with new variable $y$ and monic in $x$. We rewrite $p'''$ as a polynomial in $x$ and $y$ with coefficients in the ring $K=\mathbb F[\mathbf z]$ and call the representation $f$. That is,$f(x,y)\in K[x,y]$. By Lemma 8, the tranformation maintains irreducible factors. Note also that by the definition of $p'''$ we have $f(x,0)=p'''(x,0,0,\ldots,0)=f(x,y)\bmod y$, so that $f(x,y)\bmod y$ is univariate.
\end{enumerate}
The main part now is to factor $f(x,y)\in K[x,y]$, say $f=gh$, where $g\in K[x,y]$ is irreducible and coprime to $h\in K[x,y]$, and $f,g,h$ are monic in $x$ and have $x$-degree $\geq1$. Let $d$ be the total degree of $f$ in $x,y$.
From the factor $g$ of $f$, we will recover the factor $f_1$ of $p$ by reversing the above transformations. We show that $g$ can be computed by an ABP of size $\operatorname{poly}(s)$. It follows that the irreducible factor $f_1$ has an ABP of size $\operatorname{poly}(s)$.
The basic strategy is to first factor the univariate polynomial $f\bmod y$, and then apply Hensel lifting to get a factorization of $f\bmod y^t$, for large enough $t$. Finally, from the lifted factors modulo $y^t$, we compute the absolute factors of $f$.
\subsection*{4.1 Hensel lifting}
There are various versions of Hensel lifting in the literature (see for example~\cite{ref18}). In our case, an ABP should be able to perform several iterations of the lifting. Therefore we use the lifting in a way suitable for ABPs. In particular, in contrast to other presentations, we will not maintain the monicness of the lifted factors.
Hensel lifting works over rings $R$ modulo an ideal $I\subseteq R$. In our case,$R=K[x,y]$, where $K=\mathbb F[\mathbf z]$, and $I=\langle y\rangle^t$, for some $t\geq1$.
\begin{ntdefinition}[Author Definition 11: Lifting]
Let $R$ be a ring and $I\subseteq R$ be an ideal. Let $f,g,h,a,b\in R$ such that $f\equiv gh\pmod I$ and $ag+bh\equiv1\pmod I$. Then we call $g',h'\in R$ a lift of $g,h$ with respect to $f$, if
\begin{enumerate}
\item[(i)] $f\equiv g'h'\pmod{I^2}$,
\item[(ii)] $g'\equiv g\pmod I$ and $h'\equiv h\pmod I$, and
\item[(iii)] $\exists a',b'\in R\ a'g'+b'h'\equiv1\pmod{I^2}$.
\end{enumerate}
\end{ntdefinition}
\paragraph{Author Remark}
The three conditions in Definition 11 are the invariants when iterating the lifting.
Note that the last condition is actually redundant. It follows from the assumptions together with the second condition. This can be seen in the proof of Lemma 12 below, where a lift $g',h'$ from $g,h$ is constructed, together with $a',b'$. When we show that condition(iii) holds, we do not use the specific form of $g',h'$ constructed there, it suffices to have condition(ii).
\begin{ntlemma}[Author Lemma 12: Hensel Lifting]
Let $R$ be a ring and $I\subseteq R$ be an ideal. Let $f,g,h,a,b\in R$ such that $f\equiv gh\pmod I$ and $ag+bh\equiv1\pmod I$. Then we have:
\begin{enumerate}
\item (Existence). There exists a lift $g',h'$ of $g,h$ w.r.t.$f$.
\item (Uniqueness). For any other lift $g^*,h^*$ of $g,h$ w.r.t.$f$, there exists a $u\in I$ such that
\[g^*\equiv g'(1+u)\pmod{I^2}\quad\text{and}\quad h^*\equiv h'(1-u)\pmod{I^2}.\]
\end{enumerate}
\end{ntlemma}
\begin{proof}
We first show the existence part. Let
\begin{enumerate}\item $e=f-gh$,\item $g'=g+be$ and $h'=h+ae$,\item $c=ag'+bh'-1$,\item $a'=a(1-c)$ and $b'=b(1-c)$.\end{enumerate}
We verify that $g',h'$ is a lift of $g,h$. Because $f\equiv gh\pmod I$, we have $e=f-gh\equiv0\pmod I$. In other words,$e\in I$. It follows that $g'\equiv g\pmod I$ and $h'\equiv h\pmod I$.
Next we show that $f\equiv g'h'\pmod{I^2}$.
\begin{align*}f-g'h'&=f-(g+be)(h+ae)\\&=f-gh-e(ag+bh)-abe^2\\&\equiv e-e(ag+bh)\pmod{I^2}\\&\equiv e(1-(ag+bh))\pmod{I^2}\\&\equiv0\pmod{I^2}.\end{align*}
In the second line, note that $e^2\in I^2$. The last equality holds because $e\in I$ and $1-(ag+bh)\in I$.
Now, we verify that $a'g'+b'h'\equiv1\pmod{I^2}$. First, observe that
\begin{align*}c&=ag'+bh'-1\\&\equiv ag+bh-1\pmod I\\&\equiv0\pmod I.\end{align*}
Hence,$c\in I$ and we conclude that $a'\equiv a\pmod I$ and $b'\equiv b\pmod I$. Now,
\begin{align*}a'g'+b'h'-1&=a(1-c)g'+b(1-c)h'-1\\&=ag'+bh'-1-c(ag'+bh')\\&=c-c(ag'+bh')\\&=c(1-(ag'+bh'))\\&\equiv0\pmod{I^2}.\end{align*}
The last equality holds because $c\in I$ and $1-(ag'+bh')\equiv-c\equiv0\pmod I$.
For the uniqueness part, let $g^*,h^*$ be another lift of $g,h$. Let $\alpha=g^*-g'$ and $\beta=h^*-h'$. By Definition 11(ii), we have $g'\equiv g\equiv g^*\pmod I$ and $h'\equiv h\equiv h^*\pmod I$, and therefore $\alpha,\beta\in I$.
We first show
\begin{equation}\beta g'+\alpha h'\equiv0\pmod{I^2}.\tag{4}\end{equation}
\begin{align*}\beta g'+\alpha h'&=\beta g'+(g^*-g')h'\\&=\beta g'+g^*h'-g'h'\\&\equiv\beta g'+g^*h'-g^*h^*\pmod{I^2}\\&\equiv\beta g'-\beta g^*\pmod{I^2}\\&\equiv-\alpha\beta\pmod{I^2}\\&\equiv0\pmod{I^2}.\end{align*}
Define $u=a'\alpha-b'\beta$. Because $\alpha,\beta\in I$, also $u\in I$. Then, by(4) and because $a'g'+b'h'\equiv1\pmod{I^2}$, we have
\begin{align*}g'(1+u)&=g'(1+(a'\alpha-b'\beta))\\&=g'+a'g'\alpha-b'g'\beta\\&\equiv g'+a'g'\alpha+b'h'\alpha\pmod{I^2}\\&\equiv g'+\alpha\pmod{I^2}\\&\equiv g^*\pmod{I^2}.\end{align*}
Similarly, we get $h^*\equiv h'(1-u)\pmod{I^2}$.
\end{proof}
For the ABP-size, recall that the size just adds up when doing additions or multiplications. Hence, when $f,g,h,a,b$ have ABPs of size $\leq s$ and we construct ABPs for $g',h',a',b'$ according to steps 1--4 in the above proof, then we get ABPs of size $O(s)$.
\paragraph{Author Remark}
In the monic version of Hensel Lifting there is a division in addition to the4 steps from above. When we assume that $g$ is monic, we can compute polynomials $q$ and $r$ such that $g'-g=qg+r$, where $\deg_x(r)<\deg_x(g)$. Then one can show that $\widehat g=g+r$ and $\widehat h=h'(1+q)$ are a lift of $g,h$ w.r.t.$f$. Moreover, when $g,h$ are monic, so are $\widehat g,\widehat h$. Also the B\'ezout-coefficients $\widehat a,\widehat b$ can be computed. For $\widehat c=a\widehat g+b\widehat h-1$, let $\widehat a=a(1-\widehat c)$ and $\widehat b=b(1-\widehat c)$.
The advantage of the monic version is that the result is really unique. There is no $1+u$ factor between monic lifts. The disadvantage is the extra division which would blow up the ABP-size too much.
